\chapter{紧致空间}
\label{chap:compact-topological-spaces-solutions}

\begin{solution}{compact-fip-de-morgan}
设 \(\mathcal F\) 为闭集族，令 \(\mathcal U=\{X\setminus F:F\in\mathcal F\}\)。则
\[
 \bigcap_{F\in\mathcal F}F=\varnothing
 \Longleftrightarrow
 \bigcup_{U\in\mathcal U}U=X.
\]
而有限子族 \(F_1,\ldots,F_n\) 的交为空，恰好等价于相应有限多个补集覆盖 \(X\)。所以“每个开覆盖有有限子覆盖”与“每个总交为空的闭集族已有有限子族交为空”等价，后者的逆否命题就是有限交性质版本。空闭集族的交为 \(X\)，故在非空空间中自动非空；空空间本身紧，且空开覆盖已经是有限子覆盖，约定彼此一致。
\end{solution}

\begin{solution}{compact-subspace-criterion}
若 \(K\) 在子空间拓扑中紧，\(X\) 中开集覆盖 \(K\) 当然也是相对开覆盖，故有有限子覆盖。反之，给 \(K\) 的相对开覆盖 \(K\cap U_\alpha\)，其中 \(U_\alpha\) 在 \(X\) 中开；\((U_\alpha)\) 覆盖 \(K\)，按假设取有限子族，再与 \(K\) 相交即得所需有限相对开覆盖。
\end{solution}

\begin{solution}{compact-finite-union}
每个紧子集 \(K_j\) 从给定开覆盖中抽出有限子覆盖；有限多个有限族之并仍有限，覆盖 \(\bigcup_{j=1}^mK_j\)。

在离散空间 \(\N\) 中，每个单点集紧，但
\(\N=\bigcup_{n\in\N}\{n\}\) 不紧，因为单点开覆盖没有有限子覆盖。
\end{solution}

\begin{solution}{compact-hausdorff-bijection}
\(f(X)\) 是 \(Y\) 的子空间，仍 Hausdorff。映射 \(f:X\to f(X)\) 是连续双射；由紧空间到 Hausdorff 空间的连续双射定理，它是同胚。
\end{solution}

\begin{solution}{compact-graph}
若 \((x_0,y_0)\) 不在图像上，则 \(y_0\ne f(x_0)\)。取 \(Y\) 中不相交开邻域 \(V\ni y_0\)、\(W\ni f(x_0)\)。由 \(f\) 连续，\(U=f^{-1}(W)\) 是 \(x_0\) 的开邻域，且 \(U\times V\) 不交图像。因此图像的补集开，图像闭。

这个证明只需要 \(Y\) Hausdorff，不需要 \(X\) 紧。虽然题设还保证图像是 \(x\mapsto(x,f(x))\) 的紧像，却不能仅以此推出它闭：题设未给 \(X\) Hausdorff，积空间也未必 Hausdorff。
\end{solution}

\begin{solution}{alexander-closed-subbasis}
开子基 \(\mathcal S\) 与闭子基
\(\mathcal C=\{X\setminus S:S\in\mathcal S\}\) 一一对应。一个 \(\mathcal S\)-子族覆盖 \(X\)，当且仅当相应 \(\mathcal C\)-子族总交为空；它有有限子覆盖，当且仅当已有有限子族交为空。因此“每个子基本开覆盖有有限子覆盖”等价于“闭子基中每个有限交性质子族总交非空”，再用 Alexander 定理。
\end{solution}

\begin{solution}{maximal-fip-prime}
若 \(A\notin\mathcal M\)，则 \(\mathcal M\cup\{A\}\) 不再有有限交性质，所以存在 \(M_A\)，它是 \(\mathcal M\) 中有限多个成员之交，且 \(M_A\cap A=\varnothing\)。若同时 \(B\notin\mathcal M\)，同理得 \(M_B\cap B=\varnothing\)。于是
\[
 (M_A\cap M_B)\cap(A\cup B)=\varnothing.
\]
而 \(M_A\cap M_B\in\mathcal M\)，故 \(A\cup B\) 不能属于 \(\mathcal M\)。逆否命题给出二元结论；对有限并作归纳即可。
\end{solution}

\begin{solution}{tychonoff-finite-direct}
若因子为空，结论直接成立。以下设两因子非空，\(\mathcal U\) 为积的开覆盖。固定 \(x\in X\)，考虑所有满足
\[
 x\in V,\qquad V\times W\subseteq U\in\mathcal U
\]
的开矩形，其中 \(V\subseteq X\)、\(W\subseteq Y\)。所有这样的 \(W\) 覆盖 \(Y\)，不必先为每个 \(y\) 选定一个矩形。由紧致性，有限个 \(W_1,\ldots,W_m\) 已覆盖 \(Y\)；为它们各取对应的 \(V_k,U_k\)，令 \(O=\bigcap_{k=1}^m V_k\)。则 \(x\in O\)，且 \(O\times Y\) 被 \(U_1,\ldots,U_m\) 覆盖。

再令 \(\mathcal O\) 为所有“条带可被原覆盖有限覆盖”的开集 \(O\subseteq X\)。刚才的论证表明 \(\mathcal O\) 覆盖 \(X\)，但没有要求为每个 \(x\) 选定一个 \(O\)。紧致性给出有限子覆盖 \(O_1,\ldots,O_r\)，为这有限个条带各选一组原覆盖成员，合并即得 \(X\times Y\) 的有限子覆盖。整个证明仅作有限次选择，这可在 ZF 中由有限归纳完成。
\end{solution}

\begin{solution}{product-closed-sets}
由选择公理，从每个非空 \(F_i\) 选 \(x_i\in F_i\)，故积非空。每个
\(\pi_i^{-1}(F_i)\) 闭，所以
\[
 \prod_iF_i=\bigcap_i\pi_i^{-1}(F_i)
\]
在积空间中闭。Tychonoff 定理说明 \(\prod_iX_i\) 紧，故该闭子集紧。也可先对每个 \(F_i\) 使用闭子集紧致，再直接应用 Tychonoff。
\end{solution}
